\documentclass[11pt, oneside, a4paper]{article}
%\usepackage[cp1251]{inputenc} % кодировка
\usepackage[utf8]{inputenc} % кодировка
\usepackage[english, russian]{babel} % Русские и английские переносы
\usepackage{graphicx}          % для включения графических изображений
\usepackage{cite}              % для корректного оформления литературы
\usepackage{enumitem}
\usepackage{amsmath,amsthm,amssymb}
\usepackage{mathtext}


%стилевой пакет
\usepackage{conf2026}


\begin{document}
% \udk     - универсальный десятичный классификатор
% \msc     -
% \title   - название статьи
% \authors - список авторов

\setcounter{page}{1}

\udk{519.63}

\title{Исследование схем редукции размерности
для алгоритма поиска нуля неотрицательной функции}



\authors{А.С. Зайцев \superscript{1}, К.А. Баркалов \superscript{1}}
\organizations{Нижегородский государственный университет им. Н.И. Лобачевского
Научно-образовательный математический центр «Математика технологий будущего»\superscript{1}}



\abstractru{В работе рассматривается задача поиска глобального минимума многомерной многоэкстремальной функции, принимающей неотрицательные значения и достигающей в области определения нулевого значения. Такую задачу можно интерпретировать как задачу минимизации невязки между теоретическими и экспериментальными данными. Для ускорения сходимости алгоритма поиска нуля неотрицательной многомерной функции предлагается редуцировать размерность задачи с помощью аппроксимации кривой Пеано – неинъективной развертки. Проведены вычислительные эксперименты, подтверждающего эффективность предложенного подхода к редукции размерности в данной задаче.
}

\keywords{глобальная оптимизация, редукция размерности, кривые Пеано, поиск нуля функции.}


% \section{название} - заголовок раздела первого уровня
% \subsection{название} - заголовок раздела второго уровня
% \subsubsection{название} - заголовок раздела третьего уровня
% Не используйте уровень вложенности заголовков больше трех!
% Каждый абзац текста в статье начинается командой \par или пустой
% строкой.

\bigskip

\section{Введение}

Данная работа является продолжением исследования Алгоритма Поиска Нуля (АПН) неотрицательной функции [5, 6]. Для редукции размерности АПН использует кривые Пеано [1, 2, 3], которые отображают $N$-мерную область $D$ на единичный отрезок вещественной оси. Новым элементом исследования является использование аппроксимации кривой Пеано [4], которая для любой точки из $N$-мерной области $D$ ставит в соответствие множество точек на единичном отрезке вещественной оси. Это позволяет ускорить сходимость алгоритма к глобальному минимуму.

В проведенных вычислительных экспериментах АПН, использующий разные аппроксимации кривых Пеано, сравнивается с широко известным методом глобальной оптимизации DIRECT [7].



\section{Постановка задачи}
Классическая задача поиска минимума многомерной многоэкстремальной функции $\phi(y)$, где $y\in R^{n} (n>1)$ в $N$-мерном гиперкубе $D$ записывается в следующем виде:
	\[
	\phi(y^{*})=min\{\phi(y),y\in D\}, D=\{y\in R^{n},a_{i}\leq y_{i}\leq b_{i}, i=1,...,n\} \tag{1}
	\]

Дополнительно потребуем, чтобы функция $\phi(y)$ удовлетворяла следующим условиям:
\begin{enumerate}
	\item $\phi(y)\geq 0, y\in D$ (неотрицательность функции);
  \item Существует вектор $y^{*}\in D$ такой, что $\phi(y^{*})=0$ (существование нуля функции);
	\item Функция удовлетворяет условию: \[\left|\phi(y^{'}) - \phi(y^{''}) \right| \leq L \left| y^{'} - y^{''} \right| ,y^{'},y^{''}\in D, 0<L<\infty\]
\end{enumerate}

где значение $L$ – не задано (липшицевость функции).

Функцию $\phi(y)$, удовлетворяющую описанным выше условиям, будем обозначать $\psi(y)$, где $y\in R^{n}\:(n>1)$. Очевидно, что для рассматриваемой функции $\psi(y)\:(y\in D)$ задача поиска глобального минимума эквивалентна задачи поиска нуля функции (поиска корня функции).

Чтобы избежать экспоненциального роста вычислений при росте размерности задачи $N$ поставленная многомерная задача редуцируется до одномерной. Для редукции размерности можно использовать кривые Пеано [1, 2, 3]. Аппроксимации таких кривых (развертки) обычно обладают свойством однозначности (инъективности): любой точке (прообразу) отрезка $[0,1]$ вещественной оси соответствует единственная точка (образ) в $N$-мерной области $D$. В результате точкам, близким в области $D$, могут соответствовать расположенные далеко друг от друга  прообразы на отрезке $[0,1]$ вещественной оси. Данный эффект в некоторых случаях может негативно сказать на скорости сходимости метода оптимизации. Чтобы избежать подобного эффекта, предлагается использовать так называемую неинъективную развертку. В работе [4] предложены развертки, обладающие следующим свойством: любой точке в $N$-мерной области $D$ соответствует множество (но не более чем $2^{N}$) прообразов на отрезке $[0,1]$ вещественной оси. Таким образом исходная многомерная задача (1) сводится к одномерной задаче следующего вида:

\[
\psi(y^{*}(x))=min\{\psi(y(x)),x\in [0,1]\}\tag{2}
\]

Функция $\psi(y(x))$ удовлетворяет следующим условиям:

\begin{enumerate}
	\item $\psi(y(x))\geq 0, x\in [0,1]$ (неотрицательность функции);
  \item Существует точка $x^{*}\in [0,1]$ такая, что $\psi(y^{*}(x))=0$ (существование нуля функции);
	\item Функция удовлетворяет условию: \[\left|\psi(y(x^{'})) - \psi(y(x^{''})) \right| \leq H \sqrt[N]{\left| x^{'} - x^{''} \right|} ,x^{'},x^{''}\in [0,1], 0<H<\infty\]
\end{enumerate}
где значение $H$ – не задано. Из [2] известно, что константа Гёльдера $H$ связана с константой Липшица $L$ соотношением $H=4L\sqrt{N}$ (гёльдеровость функции).

Для решения редуцированной задачи применим Алгоритм Поиска Нуля (АПН) [5, 6]. Данный алгоритм специально разрабатывается для такого класса задач.

\section{Описание алгоритма}
Алгоритм поиска нуля функции $\psi(y(x))$ порождает последовательность точек ${x^{k}}$. В этих точках вычисляется значение функции $z^{k}=\psi(y(x^{k})), k= 0, 1, 2,...$ . Предполагается, что погрешность пренебрежимо мала. Первые два испытания осуществляются в граничных точках, то есть $x^{0}=0$ и $x^{1}=1$ соответственно. А выбор точек $x^{k+1}, k\geq 1$ , очередного $(k+1)$-го испытания осуществляются по следующим правилам:
\begin{enumerate}
	\item Перенумеруем (нижнем индексом) точки $x_{i} , 0\leq i\leq k$, в порядке возрастания значения координаты:\[0=x_{0}<x_{1}<...<x_{k}=1 \tag{3}\] 
  \item Для каждого интервала $(x_{i}-x_{i-1}), 1\leq i\leq k$, вычислим величину:\[K(i)=\frac{z_{i}+z_{i-1}}{\Delta_{i}} \tag{4}\] где $\Delta_{i}=\sqrt[N]{x_{i}-x_{i-1}},z_{i}=\psi(y(x_{i}))$, называемую характеристикой этого интервала
	\item Определим интервал $(x_{t}-x_{t-1})$, которому соответствует минимальная характеристика $K$:\[K(t)=min\{K(i),1\leq i\leq k\} \tag{5}\]Если минимальное значение соответствует нескольким интервалам, то в качестве значения $t$ выбирается наименьшее число, удовлетворяющее условию (3).
	\item Определим точку $x^{k+1}$ очередного испытания следующим образом:\[x^{k+1}=\frac{x_{t}+x_{t-1}}{2}\tag{6}\]
	\item Проведем очередное испытание в точке $x^{k+1}$:\[z^{k+1}=\psi(y(x^{k+1}))\tag{7}\] и примем $k= k+1$
	\item Проверим условие остановки алгоритма:
	\begin{description}
		\item[a.] Число итераций $k$ достигло значения $k_{max}$ (исчерпание вычислительного ресурса)
		\item[b.] Было достигнуто значение:\[\psi^{*}_{k}=min\{\psi(y(x_{i})):0\leq i\leq k\}\leq \sigma \tag{8}\]где $\sigma>0$ (предустановленная допустимая погрешность в оценке нулевого значения функции $\psi(y(x))$ )
	\end{description}
\end{enumerate}

\section{Результаты экспериментов}

В вычислительных экспериментах было проведено сравнение трех алгоритмов: АПН с кусочно-линейной разверткой, АПН с неинъективной разверткой, и метода DIRECT [7]. Во всех методах в критерии остановки использовалось значение $\sigma=10^{-2}$; в АПН использовалась развертка с параметром построения $m=10$, а в методе DIRECT – параметр $\epsilon=10^{-4}$.
Сравнение проводилось путем решения 100 задач, сформированных на основе функций вида:
\[f(y)=\sum^{14}_{i=1}(A_{i}(2i\pi y_{1})+B_{i}(2i\pi y_{2})),y\in [0,1]^{2} \tag{9}\] где значения $A_{i},B_{i}, 1\leq i\leq 14$, независимо и равномерно распределены на отрезке $[-1,1]$. Указанные функции преобразовывались к неотрицательным следующим образом:\[\phi(y)=f(y)-f^{*} \tag{10}\]где $f^{*}$ являлось известным значением в точке глобального минимума. Линии уровня одной из функций $\phi(y)$ представлены на Рис. 1; на нем же отмечены точки 55 поисковых испытаний, выполненных АПН с неинъективной разверткой (a), и точки 232 поисковых испытаний, выполненных при решении той же задачи методом DIRECT (b).

%тут должны быть картинки.
Рис. 1. График целевой функции и точек испытаний АПН с неинъективной разверткой (a) и DIRECT (b)

Результаты решения всей серии из 100 задач отражены в таблице 1, в которой приведено среднее, медианное и максимальное число испытаний, выполненных тремя методами.
%тут должна быть таблица. нужно ещё придумать таблице имя.
Табл. 1.

Результаты показывают, что АПН с неинъективной разверткой требует значительно меньшего числа поисковых испытаний как по сравнению со своим прототипом, так и по сравнению с методом DIRECT.

\section{Заключение}
По результатам проведенного вычислительного эксперимента несложно заметить, что тип развертки оказывает влияние на работу АПН. В большей степени влияние на локальную фазу поиска, т. е. когда алгоритм попал в область притяжения глобального минимума. Это связано с тем, что кусочно-линейная развертка гарантирующая однозначность  соответствия любого прообраза из единичного отрезка вещественной оси и образа в $N$-мерной области $D$ приводит к порождению минимумов, которые существуют только в редуцированной задаче, в области притяжения глобального минимума. Неинъективная развертка лишена подобного эффекта, что позволяет АПН за значительно меньшее число итераций по сравнению с методом DIRECT и по сравнению со своим прототипом решить серию из 100 задач вида (9), преобразованными к неотрицательному виду с помощью (10).

\begin{biblio}

\bibitem{Strongin1}
Стронгин Р.Г. Численные методы в многоэкстремальных задачах (информационно статистический подход) // М.: Наука, 1978, 240 стр.

\bibitem{Strongin2}
Стронгин Р.Г. Поиск глобального оптимума // М.: Знание, 1990.

\bibitem{Barkalov1}
Barkalov K. A., Strongin R. G., Bevzyuk S. A. Acceleration of Global Search by Implementing Dual Estimate for Lipschitz Constant // Lecture Notes in Computer Science. vol. 11974. 2020. P. 478-486.

\bibitem{Sergeyev}Sergeyev D. Y., Strongin R. G., Lera D. Introduction to Global Optimization Exploiting Space-Filling Curves // Springer Briefs in Optimization, 2013

\bibitem{Zaitsev}
Зайцев А. С., Баркалов К. А. Об алгоритме поиска глобального минимума для одного класса многоэкстремальных задач // Сборник трудов IV Молодёжной школы. Н.Новгород, Издательство ННГУ им. Н. И. Лобачевского, 2025. С. 39-42

\bibitem{Barkalov2}
Баркалов К. А., Зайцев А. С., Стронгин Р. Г. Алгоритм поиска нуля неотрицательной многоэкстремальной функции // Журнал вычислительной математики и математической физики. 2026. T. 66, №3. C. 348-360.

\bibitem{Jones}
Jones D. R., Perttunen C. D., Stuckman B. E. Lipschitzian optimization without the lipschitz constant // Journal of optimization theory and application: Vol. 79, No. 1, october 1993

\end{biblio}

\newpage

\msc{код MSC2010}

\title{Name of article}


\authors{R.V.~Zhalnin \superscript{1},  E.E.~Peskova \superscript{1}, V.F.~Tishkin \superscript{2}}
\organizations{Organization1 \superscript{1}, 	Organization2 \superscript{2}}

% abstract is contained in the  abstract environment
\abstracten{ Abstract.  Abstract.  Abstract. Abstract. Abstract. Abstract. Abstract. Abstract. Abstract. Abstract.
 Abstract. Abstract. Abstract. Abstract. Abstract. Abstract. Abstract. Abstract. Abstract. Abstract. Abstract. Abstract.
  Abstract. Abstract. Abstract. Abstract. Abstract. Abstract. Abstract. Abstract. Abstract. Abstract. Abstract. Abstract.
}

\keywordsen{key words, key words,key words,key words,key words,key words,key words,key words,key words,key words.}


\end{document}
